\chapter{几何光学的基本原理}

\problem{3.1}{证明反射定律符合费马原理。}
\answerarea

\sproblem{3.2}{根据费马原理可以导出在近轴光线条件下，从物点发出并会聚到像点的所有光线的光程都相等，由此导出薄透镜的物像公式。}
\answerarea
\clearpage

\problem{3.3}{眼睛 $\mathrm{E}$ 和物体 $PQ$ 之间有一块折射率为 $1.5$ 的玻璃平板（见题 3.3 图），平板的厚度 $d$ 为 $30\,\mathrm{cm}$。求物体 $PQ$ 的像 $P'Q'$ 与物体 $PQ$ 之间的距离 $d_2$ 为多少？\fig[0.52]{3_3}}
\answerarea

\problem{3.4}{玻璃棱镜的折射棱角 $A$ 为 $60^\circ$，对某一波长的光其折射率 $n$ 为 $1.6$。计算：（1）最小偏向角；（2）此时的入射角；（3）能使光线从 $A$ 角两侧透过棱镜的最小入射角。}
\answerarea
\clearpage

\problem{3.5}{题 3.5 图表示一种恒偏向棱镜，它相当于一个 $30^\circ\text{--}60^\circ\text{--}90^\circ$ 棱镜与一个 $45^\circ\text{--}45^\circ\text{--}90^\circ$ 棱镜按图示方式组合在一起。白光沿 $i$ 方向入射，旋转棱镜改变 $\theta_1$，从而使任意一种波长的光可以依次循着图示的路径传播，出射光线为 $r$。求证：如果 $\sin\theta_1=\dfrac{n}{2}$，则 $\theta_2=\theta_1$，且光束 $i$ 与 $r$ 相互垂直（这就是恒偏向棱镜名字的由来）。\fig[0.48]{3_5}}
\answerarea

\problem{3.6}{高 $5\,\mathrm{cm}$ 的物体距凹面镜顶点 $12\,\mathrm{cm}$，凹面镜的焦距是 $10\,\mathrm{cm}$，求像的位置及高度，并作光路图。}
\answerarea
\clearpage

\problem{3.7}{一个 $5\,\mathrm{cm}$ 高的物体放在球面镜前 $10\,\mathrm{cm}$ 处成 $1\,\mathrm{cm}$ 高的虚像。求：（1）此镜的曲率半径；（2）此镜是凸面镜还是凹面镜？}
\answerarea

\problem{3.8}{某观察者通过一块薄玻璃板去看在凸面镜中自己的像。他移动着玻璃板，使得在玻璃板中与在凸面镜中所看到的他眼睛的像重合在一起。若凸面镜的焦距为 $10\,\mathrm{cm}$，眼睛距凸面镜顶点的距离为 $40\,\mathrm{cm}$，问玻璃板与观察者眼睛之间的距离为多少？}
\answerarea
\clearpage

\problem{3.9}{物体位于凹面镜轴线上焦点之外，在焦点与凹镜面之间放一个与轴线垂直的两表面互相平行的玻璃板，其厚度为 $d$，折射率为 $n$。试证明：放入该玻璃板后使像移动的距离与把凹面镜向物体移动 $d(n-1)/n$ 的一段距离的效果相同。}
\answerarea

\problem{3.10}{欲使由无穷远发出的近轴光线通过透明球体并成像在右半球面的顶点处，问这透明球体的折射率应为多少？}
\answerarea
\clearpage

\problem{3.11}{有一折射率为 $1.5$、半径为 $4\,\mathrm{cm}$ 的玻璃球，物体在距球表面 $6\,\mathrm{cm}$ 处，求：（1）物所成的像到球心之间的距离；（2）像的横向放大率。}
\answerarea

\problem{3.12}{一个折射率为 $1.53$、直径为 $20\,\mathrm{cm}$ 的玻璃球内有两个小气泡，看上去一个恰好在球心，另一个从离观察者最近的方向看去，好像在表面与球心连线的中点。求两气泡的实际位置。}
\answerarea
\clearpage

\problem{3.13}{直径为 $1\,\mathrm{m}$ 的球形鱼缸的中心处有一条小鱼，若玻璃缸壁的影响可忽略不计，求缸外观察者所看到的小鱼的表观位置和横向放大率。}
\answerarea

\problem{3.14}{玻璃棒一端成半球形，其曲率半径为 $2\,\mathrm{cm}$。将它水平地浸入折射率为 $1.33$ 的水中，沿着棒的轴线离球面顶点 $8\,\mathrm{cm}$ 处的水中有一物体，利用计算和作图法求像的位置及横向放大率，并作光路图。}
\answerarea
\clearpage

\problem{3.15}{有两块玻璃薄透镜的两表面均各为凸球面及凹球面，其曲率半径为 $10\,\mathrm{cm}$。一物点在主轴上距镜 $20\,\mathrm{cm}$ 处，若物和镜均浸在水中，分别用作图法和计算法求像点的位置。设玻璃的折射率为 $1.5$，水的折射率为 $1.33$。}
\answerarea

\problem{3.16}{一凸透镜在空气中的焦距为 $40\,\mathrm{cm}$，在水中时焦距为 $136.8\,\mathrm{cm}$，问此透镜的折射率为多少（水的折射率为 $1.33$）？若将此透镜置于 $\mathrm{CS_2}$ 中（$\mathrm{CS_2}$ 的折射率为 $1.62$），其焦距又为多少？}
\answerarea
\clearpage

\problem{3.17}{两片极薄的表面玻璃，曲率半径分别为 $20\,\mathrm{cm}$ 和 $25\,\mathrm{cm}$。将两片的边缘粘起来，形成内含空气的双凸透镜，把它置于水中，求其焦距为多少？}
\answerarea

\problem{3.18}{会聚透镜和发散透镜的焦距都是 $10\,\mathrm{cm}$，求：（1）与主轴成 $30^\circ$ 的一束平行光入射到每个透镜上，像点在何处？（2）在每个透镜左方的焦平面上离主轴 $1\,\mathrm{cm}$ 处各置一发光点，成像在何处？作出光路图。}
\answerarea
\clearpage

\problem{3.19}{题 3.19 图（a）、（b）所示的 $MM'$ 分别为一薄透镜的主光轴，$S$ 为光源，$S'$ 为像。用作图法求透镜中心和透镜焦点的位置。\fig[0.64]{3_19}}
\answerarea

\sproblem{3.20}{比累对切透镜是把一块凸透镜沿直径方向剖成两半组成，两半块透镜垂直光轴拉开一点距离，用挡光的光阑 $\mathrm{K}$ 挡住其间的空隙（见题 3.20 图），这时可在屏上观察到干涉条纹。已知点光源 $P$ 与透镜相距 $300\,\mathrm{cm}$，透镜的焦距 $f'=50\,\mathrm{cm}$，两半透镜拉开的距离 $t=1\,\mathrm{mm}$，光屏与透镜相距 $l=450\,\mathrm{cm}$。用波长为 $632.8\,\mathrm{nm}$ 的氦氖激光作为光源，求干涉条纹的间距。\fig[0.52]{3_20}}
\answerarea
\clearpage

\sproblemfig[0.38]{3.21}{把焦距为 $10\,\mathrm{cm}$ 的会聚透镜的中央部分 C 切去，C 的宽度为 $1\,\mathrm{cm}$，把余下的两部分粘起来（见题 3.21 图）。如在其对称轴上距透镜 $5\,\mathrm{cm}$ 处置一点光源，试求像的位置。}{3_21}
\answerarea

\sproblemfig[0.18]{3.22}{一折射率为 $1.5$ 的薄透镜，其凸面的曲率半径为 $5\,\mathrm{cm}$，凹面的曲率半径为 $15\,\mathrm{cm}$，且镀上银（见题 3.22 图）。试证明：当光从凸表面入射时，该透镜的作用相当于一个平面镜。（提示：物经过凸面折射、凹面反射和凸面再次折射后，$s'=-s$，$\beta=1$。）}{3_22}
\answerarea
\clearpage

\problem{3.23}{题 3.23 图所示的是一个等边直角棱镜和两个透镜所组成的光学系统，棱镜折射率为 $1.5$，凸透镜的焦距为 $20\,\mathrm{cm}$，凹透镜的焦距为 $10\,\mathrm{cm}$，两透镜间距为 $5\,\mathrm{cm}$，凸透镜距棱镜边的距离为 $10\,\mathrm{cm}$。求图中长度为 $1\,\mathrm{cm}$ 的物体所成像的位置和大小。（提示：物经棱镜成像在透镜轴上，相当于经过一块厚 $6\,\mathrm{cm}$ 的平板玻璃，可利用例 3.1 的结果求棱镜所成像的位置。）\fig[0.50]{3_23}}
\answerarea

\problem{3.24}{显微镜由焦距为 $1\,\mathrm{cm}$ 的物镜和焦距为 $3\,\mathrm{cm}$ 的目镜组成，物镜与目镜之间的距离为 $20\,\mathrm{cm}$，问物体放在何处时才能使最后的像成在距离眼睛 $25\,\mathrm{cm}$ 处？作出光路图。}
\answerarea
\clearpage

\problemfig[0.40]{3.25}{题 3.25 图中 $\mathrm{L}$ 为薄透镜，水平横线 $MM'$ 为主轴，$ABC$ 为已知的一条穿过这个透镜的光线的路径，用作图法求出任一条光线 $DE$ 穿过透镜后的路径。}{3_25}
\answerarea

\sproblemfig[0.40]{3.26}{题 3.26 图中 $MM'$ 是一厚透镜的主轴，$H$、$H'$ 是透镜的主平面，$S_1$ 是点光源，$S'_1$ 是点光源的像。试用作图法求任意一物点 $S_2$ 的像 $S'_2$ 的位置。}{3_26}
\answerarea
\clearpage

\problem{3.27}{双凸薄透镜的折射率为 $1.5$，$|r_1|=10\,\mathrm{cm}$，$|r_2|=15\,\mathrm{cm}$，$r_2$ 的一面镀银，物点 $P$ 在透镜前主轴上 $20\,\mathrm{cm}$ 处，求最后像的位置并作出光路图。}
\answerarea

\problem{3.28}{实物与光屏间的距离为 $l$，在中间某一位置放一个凸透镜，可使实物的像清晰地投于屏上。将透镜移过距离 $d$ 之后，屏上又出现一个清晰的像。（1）试计算两个像的大小之比；（2）证明透镜的焦距为 $(l^2-d^2)/4l$；（3）证明 $l$ 不能小于透镜焦距的 $4$ 倍。}
\answerarea
\clearpage

\problem{3.29}{一厚透镜的焦距 $f'$ 为 $60\,\mathrm{mm}$，其两焦点间的距离为 $125\,\mathrm{mm}$，若（1）物点置于光轴上物方焦点左方 $20\,\mathrm{mm}$ 处；（2）物点置于光轴上物方焦点右方 $20\,\mathrm{mm}$ 处；（3）虚物落在光轴上像方主点右方 $20\,\mathrm{mm}$ 处，问在这三种情况下像的位置各在何处？像的性质各如何？并作光路图。}
\answerarea

\sproblem{3.30}{一个会聚薄透镜和一个发散薄透镜互相接触而成一复合光具组，当物距为 $-80\,\mathrm{cm}$ 时，实像距镜 $60\,\mathrm{cm}$，若会聚透镜的焦距为 $10\,\mathrm{cm}$，问发散透镜的焦距为多少？}
\answerarea
\clearpage

\sproblem{3.31}{双凸厚透镜两个球面表面的曲率半径分别为 $100\,\mathrm{mm}$ 和 $200\,\mathrm{mm}$，沿轴厚度为 $10\,\mathrm{mm}$，玻璃的折射率为 $1.5$，试求其焦点、主点的位置，并绘图表示之。}
\answerarea

\sproblem{3.32}{一条光线射到折射率为 $n$ 的一球形水滴（见题 3.32 图），求：（1）若后表面的入射角为 $\alpha$，问这条光线将被全反射还是部分反射？（2）偏转角 $\delta$；（3）产生最小偏转角的入射角 $\varphi$。\fig[0.52]{3_32}}
\answerarea
\clearpage

\sproblem{3.33}{将灯丝置于空心玻璃球的中心，玻璃球的内外直径分别为 $8\,\mathrm{cm}$ 和 $9\,\mathrm{cm}$。（1）求从球外观察到的灯丝像的位置（设玻璃折射率 $n=1.5$）；（2）玻璃温度计管子的内外直径分别为 $1\,\mathrm{mm}$ 和 $3\,\mathrm{mm}$，求从侧面观察到的直径的数值；（3）同一温度计竖直悬挂于直径 $100\,\mathrm{mm}$ 的盛水玻璃烧杯的正中，从较远处通过烧杯壁观察时，温度计的内、外直径分别为多少？}
\answerarea

\sproblem{3.34}{如题 3.34 图所示为梅斯林分波面干涉实验装置。其中 $O_1$、$O_2$ 分别为两块半透镜 $\mathrm{L_1}$ 和 $\mathrm{L_2}$ 的光心，$S$、$O_1$、$O_2$、$S_1$、$S_2$ 共轴，且 $S_1S_2=l$。（1）试证来自 $\mathrm{L_1}$ 和 $\mathrm{L_2}$ 两端的光束到达 $P$ 点的光程差 $\delta=l-(S_1P+S_2P)$；（2）定性讨论与轴线垂直的光屏上接收到的干涉图样的特点。\fig[0.56]{3_34}}
\answerarea
\clearpage

\sproblem{3.35}{把杂质扩散到玻璃中可以增大玻璃的折射率，这就有可能造出一个厚度均匀的透镜。已知圆板半径为 $R$，厚度为 $d$（如题 3.35 图所示），求沿半径变化的折射率 $n(r)$，它会使从 $A$ 点发出的光线传播到 $B$ 点。假定这是个薄透镜，$d\ll a$，$d\ll b$。\fig[0.52]{3_35}}
\answerarea

\sproblem{3.36}{一弯凸薄透镜两个表面的半径 $r_1$ 和 $r_2$ 分别为 $-20\,\mathrm{cm}$ 和 $-15\,\mathrm{cm}$，折射率为 $1.5$，在 $r_2$ 的凸面镀银。在距 $r_1$ 球面左侧 $40\,\mathrm{cm}$ 处的主轴上置一高为 $1\,\mathrm{cm}$ 的物，试求最后成像的位置和像的性质。}
\answerarea
\clearpage

\sproblem{3.37}{一折射率为 $n$，曲率半径为 $R_1$ 和 $R_2$ 的薄凸透镜放在折射率分别为 $n_1$ 和 $n_2$ 的两种介质之间，$s_1$ 和 $s_2$ 分别为物距和像距，$f_1$ 和 $f_2$ 是物方和像方焦距。证明：
\[\frac{f_1}{s_1}+\frac{f_2}{s_2}=1.\]}
\answerarea

\sproblem{3.38}{研究性课题：请查阅互联网，了解超薄纳米光学器件的新进展，并写成一篇小论文。}
\answerarea
\clearpage
